\chapter{Flux Superpotential and Moduli Stabilization}\label{ch:stabilization}

A compactification comes with parameters: the size of the internal space, its shape, the
value of the string coupling. In \cref{ch:k3t2} these parameters appeared as coordinates on a
moduli space, and nothing in the equations of motion preferred one point of that space to
another. In four dimensions each such parameter is a massless scalar field, and a massless
scalar that couples to matter is a long-range force that nobody has seen. This is the
\emph{moduli problem}\index{moduli problem}. The present chapter asks how a potential for the
moduli can arise, what its general shape is, and which moduli it can fix. The answer of the
last twenty-five years is: \emph{fluxes}. The quantized field strengths that \cref{ch:tadpole}
was forced to introduce in order to cancel the tadpole do not only carry charge; they also
cost energy, and the energy depends on the shape of the space through which they thread. We
shall see that this energy is summarized in one holomorphic function, the
Gukov--Vafa--Witten superpotential $W=\int G\wedge\Omega$; that the resulting potential is
non-negative because of a cancellation called \emph{no-scale structure}; that its minima are
the points where the covariant derivative $D W$ vanishes; and that one modulus, the overall
volume, always escapes. At the end we put next to this established picture the toy potential
of this programme, $V(R)\propto(R/\sqrt{\ap}-\sqrt{\ap}/R)^2$, whose minimum is the self-dual
radius of \cref{ch:tduality}. As calculus it is elementary and machine-checked; as physics it
is a proposal, and we shall be explicit about which is which.

Throughout, four-dimensional formulas are written in units where the reduced Planck mass is
one ($\kappa_4=1$), indices $I,J$ run over all chiral superfields and $i,j$ over all of them
except the volume modulus $\rho$. A bar denotes complex conjugation.

\section{The moduli problem}\label{sec:stabilization:problem}

Suppose that a compactification has a continuous family of solutions labelled by a parameter
$\phi$. Let $\phi$ vary slowly over the four non-compact dimensions. Because every constant
value of $\phi$ solves the equations of motion, the four-dimensional action for $\phi(x)$
contains no potential: it starts with the kinetic term $-\tfrac12 g(\phi)\,(\partial\phi)^2$.
Hence $\phi$ is a massless scalar field, a \emph{modulus}\index{modulus}. The radius $R$ of
the circle of \cref{ch:circle} is the simplest example; the $G_{ij}$ and $B_{ij}$ of a torus
(\cref{ch:narain}), the complex structure and K\"ahler parameters of K3 (\cref{ch:stringsk3})
and the dilaton are others.

Why is this a problem and not a feature? The argument is given by Gra\~na in three sentences
which we paraphrase. \source{Gra \S6, ll.~3914--3923}
\begin{enumerate}[leftmargin=*,itemsep=1pt]
\item A massless scalar mediates a force of infinite range, as the photon and the graviton do.
\item The coupling of a modulus to matter is not universal: the mass of the electron and the
 mass of the proton depend on the size of the internal space in different ways. Two test
 bodies of different composition would therefore fall with different accelerations, in
 conflict with tests of the equivalence principle, which hold to one part in $10^{13}$.
\item If the expectation value of the volume modulus ``can be anything'', the theory predicts
 no couplings at all, since four-dimensional couplings are functions of the moduli.
\end{enumerate}
The conclusion is that in a realistic vacuum ``all `moduli' should be massive''. A mass is the
second derivative of a potential at its minimum, so we need a potential $V(\phi)$ with an
isolated minimum. Two warnings are in order before we look for one. First, a potential
generated by a classical effect typically falls to zero at infinite volume or zero coupling,
where the effect switches off; a generic potential therefore has a \emph{runaway} direction\index{runaway direction}
and no minimum at finite distance. Gra\~na notes that ``for the GVW superpotential, the resulting
potential has flat or runaway directions''. \source{Gra \S6, ll.~3928--3930} Second, the amount
of supersymmetry matters. With sixteen unbroken supercharges, as for type II strings on $\KT$
(\cref{ch:k3t2}), the scalar manifold is a fixed symmetric space and supersymmetry allows no
potential on it (standard; not checked against a source in this book's library). A potential
requires breaking to at most four supercharges, $N=1$ in four dimensions, which is what
orientifold planes and fluxes achieve in the type IIB model $\KT/\Z_2$ of \cref{ch:tadpole}.

\section{The Gukov--Vafa--Witten superpotential}\label{sec:stabilization:gvw}

\subsection{Flux restricts the geometry}

Take M-theory on a compact Calabi--Yau fourfold $Y$ with four-form flux $G$, the setting of
\cref{ch:tadpole}.\index{Calabi--Yau fourfold} For $G=0$ the complex structure parameters $t_i$ (there are $h^{3,1}$ of
them) and the K\"ahler parameters $k_j$ (there are $h^{1,1}$) are arbitrary. For $G\neq0$ this
is no longer so. Supersymmetry requires \source{GVW \S2.4, ll.~895--903}
\begin{equation}\label{eq:stabilization:susyG}
 G^{0,4}=G^{1,3}=0,\qquad G\wedge J=0 ,
\end{equation}
where $G^{p,q}$ is the component of $G$ with $p$ holomorphic and $q$ antiholomorphic indices
and $J$ is the K\"ahler form (GVW write $K$ for it; we keep $K$ for the K\"ahler potential).
The logic of \eqref{eq:stabilization:susyG} deserves a moment of thought. The flux $G$ is an
\emph{integral} cohomology class, fixed once and for all by quantization
(\cref{ch:tadpole}); it does not know about the complex structure. The decomposition of a
four-form into types $(p,q)$, on the other hand, depends on the complex structure: when the
$t_i$ change, the same class $G$ acquires different components. Hence
\eqref{eq:stabilization:susyG} is not a condition on $G$ but a set of \emph{equations for the
moduli}, with $G$ as a parameter. This is the mechanism of flux stabilization in one sentence.

\subsection{A superpotential that reproduces the conditions}

In a theory with four supercharges the potential for chiral fields is governed by a holomorphic
function, the superpotential\index{superpotential} $W$, and in flat space the vacua preserving
supersymmetry are the solutions of $W=\dd W=0$. Gukov, Vafa and Witten looked for a $W$
whose supersymmetric vacua are \eqref{eq:stabilization:susyG} and proposed
\source{GVW \S2.4, eq.~(2.38), ll.~943--951}
\begin{equation}\label{eq:stabilization:gvw}
 W(t_i)=\frac{1}{2\pi}\int_Y\Omega\wedge G ,
\end{equation}
with $\Omega$ the holomorphic four-form of $Y$\index{Gukov--Vafa--Witten superpotential}. Let us
check the claim. The form $\Omega$ is of type $(4,0)$, and the wedge product of a $(4,0)$
form with a $(p,q)$ form can be integrated over $Y$ only if $(p,q)=(0,4)$. Therefore
\begin{equation}\label{eq:stabilization:W04}
 W=0\iff \int_Y\Omega\wedge G^{0,4}=0\iff G^{0,4}=0,
\end{equation}
the last step because $H^{0,4}(Y)$ is one-dimensional and spanned by $\bar\Omega$, with
$\int\Omega\wedge\bar\Omega\neq0$. Next, vary the complex structure. The derivative
$\partial\Omega/\partial t_i$ is no longer of pure type: it has a $(4,0)$ part proportional to
$\Omega$ and a $(3,1)$ part $\chi_i$, and the $\chi_i$ span $H^{3,1}(Y)$. Given $W=0$, the
condition $\partial_iW=0$ says that $\int\chi_i\wedge G=0$ for all $i$, that is, $G^{1,3}=0$.
``So we have found the supersymmetric interaction that accounts for'' the first two
equations in \eqref{eq:stabilization:susyG}. \source{GVW \S2.4, ll.~961--965} The third one,
$G\wedge J=0$, follows in the same way from a second function
$\widetilde W=\int_Y\mathcal K\wedge\mathcal K\wedge G$ of the complexified K\"ahler class
$\mathcal K$. \source{GVW \S2.4, eq.~(2.41), ll.~979--988}

\begin{physicsbox}[Why $W$ must be this integral: domain walls]
GVW give a second, more physical argument. A five-brane wrapped on a four-cycle $S\subset Y$
looks in the non-compact directions like a wall. On its two sides the flux differs,
$G_2-G_1=2\pi[S]$: the brane is the source of $G$, so crossing it changes the flux by one unit
through the dual cycle. The tension of the wall is the volume of $S$, which for a
supersymmetric cycle is $|\int_S\Omega|$. In a theory with four supercharges the tension of a
supersymmetric domain wall equals the absolute value of the jump of $W$ across it. Writing
$\int_S\Omega=\frac1{2\pi}\int_Y\Omega\wedge(G_2-G_1)$ one reads off $\Delta W$, ``a statement
that is clearly compatible with'' \eqref{eq:stabilization:gvw}.
\source{GVW \S2.4, eq.~(2.40), ll.~966--977} The superpotential is thus fixed by an
object whose tension can be computed independently.
\end{physicsbox}

\subsection{Type IIB: three-form flux and the axio-dilaton}

Our main example lives in type IIB string theory, where the relevant fluxes are the
Ramond--Ramond three-form $F_3$ and the Neveu--Schwarz three-form $H_3$\index{flux compactification} on a Calabi--Yau
threefold (or orientifold) $M$. They enter through one complex combination,
\source{GKP \S2.1, ll.~168--169}
\begin{equation}\label{eq:stabilization:G3}
 G_3=F_3-\tau H_3,\qquad \tau=C_0+\ii\,\ee^{-\phi},
\end{equation}
where $\tau$ is the axio-dilaton\index{axio-dilaton}: $C_0$ is the Ramond--Ramond scalar and
$\ee^{\phi}=g_s$ the string coupling, so that weak coupling means large $\operatorname{Im}\tau$.
The superpotential is \source{GKP \S2.3, eq.~(2.40), ll.~799--808; Gra \S5.4, eq.~(5.43),
ll.~3500--3509}
\begin{equation}\label{eq:stabilization:WIIB}
 W=\int_M\Omega\wedge G_3=\int_M\Omega\wedge F_3-\tau\int_M\Omega\wedge H_3 ,
\end{equation}
with $\Omega$ now the holomorphic three-form. It is the F-theory limit of
\eqref{eq:stabilization:gvw}: on an elliptically fibred fourfold a four-form flux with one
leg on the fibre decomposes as $G=\sum_iH_i\wedge\theta^i$ with two three-forms $H_i$ on the
base and $\theta^i$ the two one-forms of the fibre, \source{GVW \S2.5, eq.~(2.42),
ll.~1085--1096} and the pair $(H_1,H_2)$ corresponds to $(F_3,H_3)$; GKP give the precise
map, $G_4=-G_3\wedge\dd\bar w/(\tau-\bar\tau)+\text{h.c.}$ with $w$ the fibre coordinate.
\source{GKP \S2.3, eqs.~(2.41)--(2.42), ll.~809--840} Three features of
\eqref{eq:stabilization:WIIB} control everything that follows.
\begin{itemize}[leftmargin=*,itemsep=1pt]
\item $W$ depends on the complex structure moduli $z^\alpha$ through the periods of $\Omega$.
\item $W$ depends on $\tau$, linearly. The dilaton couples differently to the two fluxes, and
 for this reason fluxes ``generically'' stabilize it. \source{GKP \S1, ll.~97--100}
\item $W$ does not depend on the K\"ahler moduli at all: ``the K\"ahler moduli [\dots] do not
 appear in the superpotential''. \source{Gra \S5.4, ll.~3506--3509}
\end{itemize}
If the three-cycles of $M$ are $A_K,B^K$ with periods $Z^K=\int_{A_K}\Omega$,
$\mathcal F_K=\int_{B^K}\Omega$, and the fluxes through them are integers (in units of
$(2\pi)^2\ap$), then \eqref{eq:stabilization:WIIB} is a finite sum: integers times periods.
For the deformed conifold\index{conifold} with $M$ units of $F_3$ through the $A$-cycle and $-K$ units of
$H_3$ through the $B$-cycle, GKP obtain $W=(2\pi)^2\ap\,\bigl(M\,\mathcal G(z)-K\tau z\bigr)$ with
$z=\int_A\Omega$ and $\mathcal G(z)=\int_B\Omega=\frac{z}{2\pi\ii}\ln z+{}$holomorphic.
\source{GKP \S3, eqs.~(3.4)--(3.10), ll.~1007--1090} We shall use a simpler model of the same
type in \cref{sec:stabilization:worked}.

\begin{remark}[$W$ is a section, not a function]\label{rem:stabilization:section}
The form $\Omega$ is defined only up to multiplication by a nowhere-vanishing holomorphic
function of the moduli, $\Omega\mapsto\ee^{-f(z)}\Omega$. Then $W\mapsto\ee^{-f}W$: strictly
speaking $W$ is ``a section of a line bundle over the moduli space''. \source{GVW \S2.4,
ll.~952--961} An observable cannot depend on this choice, and we shall verify below that the
potential does not.
\end{remark}

\section{The \texorpdfstring{$N=1$}{N = 1} F-term potential}\label{sec:stabilization:fterm}

\subsection{Ingredients}

A four-dimensional $N=1$ supergravity theory with chiral superfields $\phi^I$ is specified, as
far as the scalar fields are concerned, by two functions: a real K\"ahler
potential\index{K\"ahler potential} $K(\phi,\bar\phi)$ and a holomorphic superpotential
$W(\phi)$. The kinetic terms are $-K_{I\bar J}\,\partial_\mu\phi^I\partial^\mu\bar\phi^{\bar J}$
with the K\"ahler metric $K_{I\bar J}=\partial_I\partial_{\bar J}K$, which must be positive
definite, and $K^{I\bar J}$ denotes its inverse. The potential is, in the absence of D-terms,
\source{Gra \S5.2, eqs.~(5.21)--(5.22), ll.~2990--3000; GKP \S2.3, eq.~(2.43), ll.~855--869}
\begin{equation}\label{eq:stabilization:V}
 V=\ee^{K}\Bigl(K^{I\bar J}\,D_IW\,\overline{D_JW}-3\,|W|^2\Bigr),\qquad
 D_IW=\partial_IW+(\partial_IK)\,W .
\end{equation}
This is the \emph{F-term potential}\index{F-term potential}; $D_IW$ is the K\"ahler covariant
derivative\index{K\"ahler covariant derivative}\index{supersymmetry breaking}, and $F_I=\ee^{K/2}D_IW$ is the order parameter
of supersymmetry breaking. GKP carry an overall factor $1/2\kappa_{10}^2$ from the
ten-dimensional normalization; it plays no role for the location of minima.

Two structural facts should be checked once by every student.

\paragraph{K\"ahler invariance.}\index{K\"ahler invariance} Under $K\mapsto K+f+\bar f$, $W\mapsto\ee^{-f}W$ with $f$
holomorphic, the kinetic terms are unchanged because $\partial_I\partial_{\bar J}(f+\bar f)=0$.
For the covariant derivative,
\[
 D_I(\ee^{-f}W)=\ee^{-f}\bigl(\partial_IW-f_IW\bigr)+(K_I+f_I)\,\ee^{-f}W=\ee^{-f}D_IW ,
\]
so $D_IW$ transforms like $W$, which the plain derivative $\partial_IW$ does not. Finally
$\ee^{K}\mapsto\ee^{K}|\ee^{f}|^2$ compensates $|\ee^{-f}|^2$ from the bracket, and $V$ is
invariant. For the flux superpotential this is exactly the freedom of
\cref{rem:stabilization:section}, because the K\"ahler potential of the complex structure
moduli is $-\ln(-\ii\int_M\Omega\wedge\bar\Omega)$, see \eqref{eq:stabilization:Ktot} below.

\paragraph{The sign structure.} The first term in \eqref{eq:stabilization:V} is a positive
definite quadratic form in the $D_IW$; the second is negative. Global supersymmetry knows only
the first. The term $-3|W|^2$ is gravitational: it is the reason why supergravity, unlike
global supersymmetry, admits supersymmetric vacua of negative energy.

\begin{proposition}[Supersymmetric points are critical points]\label{prop:stabilization:crit}
If $D_IW=0$ for all $I$ at a point $\phi_0$, then $\partial_IV(\phi_0)=0$ for all $I$, and
$V(\phi_0)=-3\,\ee^{K}|W|^2\le0$.
\end{proposition}
\begin{proof}
Write $V=\ee^K Q-3\,\ee^{K}W\bar W$ with $Q=K^{J\bar L}D_JW\overline{D_LW}$. By the product rule
every term of $\partial_I(\ee^KQ)$ contains a factor $D_JW$ or $\overline{D_LW}$ and
vanishes at $\phi_0$. For the second piece, $\bar W$ is antiholomorphic, so
$\partial_I(\ee^KW\bar W)=\ee^K\bar W(K_IW+\partial_IW)=\ee^K\bar W\,D_IW$, which vanishes as
well. The value of $V$ is read off from \eqref{eq:stabilization:V}.
\end{proof}
A supersymmetric vacuum is therefore Minkowski if $W=0$ there and anti-de Sitter if
$W\neq0$; it is never de Sitter. The proposition does not say that a supersymmetric point is
a \emph{minimum}; this needs further information, which the no-scale structure will supply.

\subsection{The K\"ahler potential of a type IIB flux compactification}

For one K\"ahler modulus $\rho$ (the overall volume; $\operatorname{Im}\rho\sim R^4$ in
string units), the dilaton and the complex structure moduli, dimensional reduction of the
ten-dimensional action gives \source{GKP \S2.3, eqs.~(2.37)--(2.38), ll.~773--789}
\begin{equation}\label{eq:stabilization:Ktot}
 K=-3\ln\bigl[-\ii(\rho-\bar\rho)\bigr]-\ln\bigl[-\ii(\tau-\bar\tau)\bigr]
 -\ln\Bigl[-\ii\int_M\Omega\wedge\bar\Omega\Bigr].
\end{equation}
Note that $-\ii(\rho-\bar\rho)=2\operatorname{Im}\rho>0$, so the logarithms are real. Gra\~na
uses $T=-\ii\rho$, in terms of which the first term reads $-3\ln(T+\bar T)$.
\source{Gra \S5.2, eq.~(5.27), ll.~3075--3083} The coefficient $3$ is the important number of
this chapter.

\section{No-scale structure}\label{sec:stabilization:noscale}

\subsection{The cancellation}

Because $W$ in \eqref{eq:stabilization:WIIB} does not depend on $\rho$, the covariant
derivative in the $\rho$ direction is purely the connection term, $D_\rho W=K_\rho W$. From
$K^{(\rho)}=-3\ln[-\ii(\rho-\bar\rho)]$,
\begin{equation}\label{eq:stabilization:Krho}
 K_\rho=-\frac{3}{\rho-\bar\rho},\qquad K_{\bar\rho}=+\frac{3}{\rho-\bar\rho},\qquad
 K_{\rho\bar\rho}=-\frac{3}{(\rho-\bar\rho)^2}=\frac{3}{4(\operatorname{Im}\rho)^2}>0 ,
\end{equation}
and therefore
\begin{equation}\label{eq:stabilization:noscale}
 K^{\rho\bar\rho}K_\rho K_{\bar\rho}
 =\Bigl(-\frac{(\rho-\bar\rho)^2}{3}\Bigr)\Bigl(-\frac{9}{(\rho-\bar\rho)^2}\Bigr)=3 .
\end{equation}
This is the \emph{no-scale identity}\index{no-scale structure}. The $\rho$ contribution to the
first term of \eqref{eq:stabilization:V} is $K^{\rho\bar\rho}|K_\rho|^2|W|^2=3|W|^2$, and it
cancels the negative term exactly:
\begin{equation}\label{eq:stabilization:Vns}
 V=\ee^{K}\,K^{i\bar\jmath}\,D_iW\,\overline{D_jW}\;\ge0 ,
\end{equation}
where $i,j$ run over $\tau$ and the $z^\alpha$ only (the metric is block diagonal because $K$ is
a sum). ``In no-scale models the $|D_\rho W|^2$ term cancels the negative term, leaving a
nonnegative potential.'' \source{GKP \S2.3, ll.~866--869} \TL{} The coefficient matters: for
$K^{(\rho)}=-p\ln[-\ii(\rho-\bar\rho)]$ the same computation gives
$K^{\rho\bar\rho}K_\rho K_{\bar\rho}=p$ and a residual term $(p-3)|W|^2$ of either sign
(symbolic check with sympy for general $p$). The value $p=3$ is tied to the fact that
$\ee^{-K^{(\rho)}}$ is proportional to the square of the six-dimensional volume in
four-dimensional Planck units.

With several K\"ahler moduli the statement is $K^{I\bar J}\partial_IK\,\partial_{\bar J}K=3$,
\source{Gra \S5.2, eq.~(5.28), ll.~3086--3096} and it has an equally short proof.

\begin{lemma}[Homogeneity implies no-scale]\label{lem:stabilization:hom}
Let $K=-\ln Y(x)$ where $Y$ is a positive function of the real variables
$x^I=T^I+\bar T^I$, homogeneous of degree $p$, and assume $K_{IJ}$ invertible. Then
$K^{IJ}K_IK_J=p$.
\end{lemma}
\begin{proof}
Since $K$ depends on $T^I$ only through $x^I$, $\partial_{T^I}\partial_{\bar T^J}K=K_{IJ}$, the
ordinary Hessian. Euler's relation $x^IY_I=pY$ gives $x^IK_I=-p$. Differentiating with respect
to $x^J$: $K_J+x^IK_{IJ}=0$, so $x^I=-K^{IJ}K_J$. Inserting this back,
$-p=x^IK_I=-K^{IJ}K_JK_I$.
\end{proof}
For a Calabi--Yau orientifold $Y$ is the square of the volume, which is homogeneous of degree
$3/2$ in the four-cycle volumes $\operatorname{Re}T^I$; hence $p=3$. We have checked the lemma
symbolically for $Y=x^1x^2x^3$ and $Y=x^1(x^2)^2+(x^3)^3$, and numerically for
$Y=\bigl((x^1)^{3/2}-(x^2)^{3/2}-(x^3)^{3/2}\bigr)^2$.

\subsection{Consequences}

Equation \eqref{eq:stabilization:Vns} has four consequences, all stated by GKP.
\source{GKP \S2.3, ll.~866--872}
\begin{enumerate}[leftmargin=*,itemsep=1pt]
\item $V\ge0$, and $V=0$ exactly where $D_iW=0$ for all $i\neq\rho$. These points are global
 minima: \emph{the minima of the flux potential are at $D_iW=0$}.
\item These are $n$ complex equations for the $n$ fields $\tau,z^\alpha$. Generically the
 solutions are isolated, and all these fields are massive.
\item The conditions do not involve $\rho$, which ``represents $n$ equations on $n$ moduli and
 leaves $\rho$ undetermined'': the volume is a flat direction with $V=0$ along it.
\item ``Generically at these solutions $W\neq0$, so $D_\rho W=-3W/(\rho-\bar\rho)$ is nonzero
 and supersymmetry is broken.''
\end{enumerate}
The last two items explain the name. Supersymmetry is broken, yet the vacuum energy is zero,
and the scale of the breaking, measured by the gravitino mass\index{gravitino mass} $m_{3/2}=\ee^{K/2}|W|\propto
(\operatorname{Im}\rho)^{-3/2}$ (standard; not checked against a source in this book's
library), slides with the undetermined field $\rho$: the theory has \emph{no scale}.

\begin{table}[t]
\centering\small
\begin{tabular}{@{}llll@{}}
\toprule
component of $G_3$ & detected by & ISD? & effect\\
\midrule
$(2,1)$, primitive & --- & yes & allowed, supersymmetric\\
$(0,3)$ & $W=\int\Omega\wedge G_3$ & yes & $V=0$, breaks SUSY via $D_\rho W$\\
$(3,0)$ & $D_\tau W\propto\int\Omega\wedge\bar G_3$ & no & $V>0$\\
$(1,2)$ & $D_\alpha W=\int\chi_\alpha\wedge G_3$ & no & $V>0$\\
\bottomrule
\end{tabular}
\caption{Hodge types of the three-form flux and the four-dimensional quantities that detect
them; ISD means imaginary self-dual, $\ast_6G_3=\ii G_3$. \source{GKP \S2.2--2.3, ll.~629--632, 706--709, 874--896; Gra \S5.4, eqs.~(5.44)--(5.47),
ll.~3520--3552}}\label{tab:stabilization:hodge}
\end{table}

\subsection{What the conditions mean in ten dimensions}

With \eqref{eq:stabilization:Ktot} one has $K_\tau=-1/(\tau-\bar\tau)$ and
$\partial_\tau W=-\int\Omega\wedge H_3$. Therefore
\begin{equation}\label{eq:stabilization:Dtau}
 D_\tau W=-\int\Omega\wedge H_3-\frac{1}{\tau-\bar\tau}\int\Omega\wedge(F_3-\tau H_3)
 =\frac{1}{\bar\tau-\tau}\int_M\Omega\wedge\bar G_3,\qquad \bar G_3=F_3-\bar\tau H_3 ,
\end{equation}
since $-(\tau-\bar\tau)H_3-(F_3-\tau H_3)=-(F_3-\bar\tau H_3)$. This is GKP's eq.~(2.44) (the
bar over $G_{(3)}$ is lost in the text extraction; Gra\~na's eq.~(5.44) writes it as
$\tfrac{\ii}{2}\ee^{\phi}\int\bar G_3\wedge\Omega$, which agrees because
$1/(\bar\tau-\tau)=\ii\,\ee^{\phi}/2$). For the complex structure moduli
$D_\alpha W=\int\chi_\alpha\wedge G_3$ with $\chi_\alpha$ a basis of $(2,1)$ forms.
\source{GKP \S2.3, eq.~(2.44), ll.~874--896} Reading off Hodge types as in
\cref{sec:stabilization:gvw}: $D_\tau W=0$ removes the $(3,0)$ part of $G_3$, $D_\alpha W=0$
removes the $(1,2)$ part, and what remains, $(2,1)\oplus(0,3)$, is precisely the flux that is
\emph{imaginary self-dual}\index{imaginary self-dual flux}, $\ast_6G_3=\ii G_3$, the condition GKP
derive from the ten-dimensional equations of motion. \source{GKP \S2.2, eq.~(2.31),
ll.~629--632, 706--709} The $(0,3)$ part is measured by $W$ itself and is the part that breaks
supersymmetry. \Cref{tab:stabilization:hodge} summarizes the dictionary. The agreement of the
four-dimensional conditions with the ten-dimensional ones is the main evidence that
\eqref{eq:stabilization:WIIB} is the correct superpotential.

\begin{pitfallbox}[``No-scale'' is a classical statement]
The flat direction and the vanishing vacuum energy hold in the classical supergravity
approximation. GKP warn that ``there is no known reason that it should survive quantum
corrections'' and that already $\ap$ corrections ``presumably spoil the no-scale structure''.
\source{GKP \S2.3, ll.~756--758} GVW make the same point for fourfolds. \source{GVW \S2.4,
ll.~1000--1008 and footnote 11, ll.~1036--1039} One should therefore not read \eqref{eq:stabilization:Vns} as
a solution of the cosmological constant problem, and not read the flat direction as a
prediction. It is the leading term of an expansion.
\end{pitfallbox}

\section{Worked example: fixing the dilaton with two periods}\label{sec:stabilization:worked}

We now solve a model completely. Freeze the complex structure, so that the periods of
$\Omega$ are fixed complex numbers, and keep two of them, $\Pi_1=1$ and $\Pi_2=\ii$. With
integer fluxes $(f_1,f_2)$ of $F_3$ and $(h_1,h_2)$ of $H_3$ through the two cycles,
\eqref{eq:stabilization:WIIB} becomes
\begin{equation}\label{eq:stabilization:toyW}
 W(\tau)=A-\tau B,\qquad A=f_1+\ii f_2,\quad B=h_1+\ii h_2 .
\end{equation}
This is a caricature: on a real Calabi--Yau the periods move with $z^\alpha$ and the fluxes are
constrained by the tadpole. It is, however, the two-period version of exactly the toy
superpotential that the Lean file of \cref{sec:stabilization:lean} defines, and it shows every
feature of the general discussion. The fields are $\tau$ and $\rho$, and
$K=-\ln(2\operatorname{Im}\tau)-3\ln(2\operatorname{Im}\rho)$.

\paragraph{Step 1: the F-term.} By \eqref{eq:stabilization:Dtau},
$D_\tau W=-(A-\bar\tau B)/(\tau-\bar\tau)$. It vanishes at
\begin{equation}\label{eq:stabilization:tau0}
 \bar\tau_0=\frac AB\quad\Longleftrightarrow\quad
 \tau_0=\frac{(f_1-\ii f_2)(h_1+\ii h_2)}{h_1^2+h_2^2},\qquad
 \operatorname{Im}\tau_0=\frac{f_1h_2-f_2h_1}{h_1^2+h_2^2}.
\end{equation}
GKP's conifold model has the same structure, $\bar\tau=M\mathcal G(0)/K'z'(0)$.
\source{GKP \S3, eqs.~(3.15)--(3.17), ll.~1139--1160} The string coupling must be positive, so
we need
\begin{equation}\label{eq:stabilization:Nflux}
 N_{\rm flux}:=f_1h_2-f_2h_1>0 .
\end{equation}
This integer is the intersection pairing of the two flux vectors: it is the toy version of
$\int H_3\wedge F_3$, the D3-brane charge carried by the flux, which entered the tadpole
condition of \cref{ch:tadpole} with a definite sign. \source{GKP \S3, eq.~(3.5),
ll.~1019--1029} The same sign that makes the flux contribute positively to the tadpole makes
the stabilized coupling positive.

\paragraph{Step 2: the potential.} With $K^{\tau\bar\tau}=-(\tau-\bar\tau)^2$ the factors of
$\tau-\bar\tau$ cancel in $K^{\tau\bar\tau}|D_\tau W|^2=|A-\bar\tau B|^2$, the no-scale identity
removes $-3|W|^2$, and $\ee^K=1/(16\operatorname{Im}\tau\,(\operatorname{Im}\rho)^3)$:
\begin{equation}\label{eq:stabilization:toyV}
 V(\tau,\rho)=\frac{|A-\bar\tau B|^2}{16\,\operatorname{Im}\tau\,(\operatorname{Im}\rho)^3}
 =\frac{|B|^2\,|\tau-\tau_0|^2}{16\,\operatorname{Im}\tau\,(\operatorname{Im}\rho)^3}\;\ge0 .
\end{equation}
(Symbolic check of the cancellation with sympy.) The potential vanishes only at $\tau=\tau_0$,
for every value of $\rho$; away from $\tau_0$ it is positive and decreases like
$(\operatorname{Im}\rho)^{-3}$, a runaway towards infinite volume.

\paragraph{Step 3: numbers.} Take $(f_1,f_2)=(2,0)$ and $(h_1,h_2)=(0,1)$. Then $A=2$,
$B=\ii$, $N_{\rm flux}=2$ and $\tau_0=2\ii$: the axion is fixed at $C_0=0$ and the string
coupling at $g_s=1/\operatorname{Im}\tau_0=\tfrac12$. Along the line $\tau=\ii t$ one finds
$A-\bar\tau B=2-t$ and $V\propto(2-t)^2/t$, shown in \cref{fig:stabilization:toy}. The value
of the superpotential at the minimum is
\begin{equation}\label{eq:stabilization:W0}
 W_0=A-\tau_0B=-\frac{2\ii N_{\rm flux}}{\bar B}=4\neq0 ,
\end{equation}
so by item 4 above supersymmetry is broken, with $D_\rho W=-3W_0/(\rho-\bar\rho)$, while
$V=0$. In this rigid toy $W_0\neq0$ is unavoidable, since \eqref{eq:stabilization:W0} vanishes
only for $N_{\rm flux}=0$; in a model with complex structure moduli one may try to tune
$W_0=0$ or $|W_0|\ll1$. A second choice, $(f_1,f_2)=(1,0)$ and $(h_1,h_2)=(1,1)$, gives
$N_{\rm flux}=1$, $\tau_0=\tfrac12(1+\ii)$ and $W_0=1-\ii$; this is strong coupling,
$g_s=2$, and one should use the $\SLZ$ duality of type IIB to map it to the fundamental domain
of \cref{ch:torus2} (\cref{ex:stabilization:sl2}).

\begin{figure}[t]
\centering
\begin{tikzpicture}[scale=0.95]
 \draw[-{Stealth}] (0,0) -- (7.4,0) node[below] {$t=\operatorname{Im}\tau$};
 \draw[-{Stealth}] (0,0) -- (0,3.6) node[left] {$V$};
 \draw[thick,domain=0.62:7,samples=80] plot (\x,{0.55*(2-\x)*(2-\x)/\x});
 \draw[dashed,domain=0.38:7,samples=80] plot (\x,{0.2*(2-\x)*(2-\x)/\x});
 \fill (2,0) circle (1.6pt) node[below] {$t_0=2$};
 \node[right] at (0.8,3.2) {\small $\operatorname{Im}\rho$ small};
 \node[right] at (4.2,0.45) {\small $\operatorname{Im}\rho$ larger};
 \draw (1,0.06)--(1,-0.06) node[below] {\small $1$};
 \draw (4,0.06)--(4,-0.06) node[below] {\small $4$};
 \draw (6,0.06)--(6,-0.06) node[below] {\small $6$};
\end{tikzpicture}
\caption{The toy flux potential \eqref{eq:stabilization:toyV} for $A=2$, $B=\ii$ along
$\tau=\ii t$: $V\propto(2-t)^2/t$. The minimum at $g_s=1/t_0=\tfrac12$ has $V=0$ for every
value of the volume modulus; increasing $\operatorname{Im}\rho$ lowers the whole curve (dashed)
without moving the minimum. The volume is not stabilized.}\label{fig:stabilization:toy}
\end{figure}

\begin{pitfallbox}[One period is not enough]
With a single period $\Pi$, $W=(f-\tau h)\Pi$ and $D_\tau W=0$ gives $\bar\tau=f/h$, a
\emph{real} number: there is no vacuum in the upper half plane. GKP meet this obstruction in the
conifold: near $z=0$ their superpotential is effectively independent of $\tau$ and ``the
$D_\tau$ equation cannot be satisfied''; the cure is ``additional fluxes'' on a second pair of
cycles. \source{GKP \S3, ll.~1118--1139} Stabilizing the dilaton needs at least two periods
whose ratio is not real. The Lean structure \lean{GVWSuperpotential} of
\cref{sec:stabilization:lean} has one period, so its toy model has no supersymmetric vacuum
with $h_3\neq0$; this is a limitation of the toy, not a statement about string theory.
\end{pitfallbox}

\section{Beyond no-scale: fixing the volume}\label{sec:stabilization:kklt}

Since $W$ receives no perturbative corrections, a dependence on the K\"ahler modulus can enter
the superpotential only non-perturbatively, from Euclidean D3-branes or gaugino condensation
on D7-branes, in the form $W_{\rm np}=B\,\ee^{-2\pi T}$ with $T=-\ii\rho$.
\source{Gra \S7.1--7.2, ll.~4738--4745, 4818--4826} Kachru, Kallosh, Linde and Trivedi (KKLT)
assume that $\tau$ and $z^\alpha$ have been fixed by fluxes\index{non-perturbative superpotential} as above, leaving a constant
$W_0$, and study\index{KKLT scenario}
\begin{equation}\label{eq:stabilization:kklt}
 W=W_0+B\,\ee^{-2\pi T},\qquad K=-3\ln(T+\bar T).
\end{equation}
\source{Gra \S7.3, eq.~(7.7), ll.~4855--4871} Set the axion to zero, $T=\sigma$ real, and take
$W_0,B$ real. Then $D_TW=-2\pi B\,\ee^{-2\pi\sigma}-\frac{3}{2\sigma}(W_0+B\,\ee^{-2\pi\sigma})$,
and $D_TW=0$ gives
\begin{equation}\label{eq:stabilization:kkltW0}
 W_0=-B\,\ee^{-2\pi\sigma_{\rm cr}}\Bigl(1+\frac{4\pi}{3}\sigma_{\rm cr}\Bigr),
\end{equation}
which is Gra\~na's eq.~(7.8). At this point $W=-\tfrac{4\pi}{3}\sigma_{\rm cr}B\,
\ee^{-2\pi\sigma_{\rm cr}}$, and by \cref{prop:stabilization:crit}
\begin{equation}\label{eq:stabilization:kkltV}
 V_{\min}=-3\,\ee^K|W|^2=-\frac{3}{8\sigma_{\rm cr}^3}\cdot\frac{16\pi^2\sigma_{\rm cr}^2B^2
 \ee^{-4\pi\sigma_{\rm cr}}}{9}=-\frac{2\pi^2B^2\,\ee^{-4\pi\sigma_{\rm cr}}}{3\,\sigma_{\rm cr}} ,
\end{equation}
in agreement with eq.~(7.9) of the same source. \source{Gra \S7.3, eqs.~(7.8)--(7.9),
ll.~4873--4898} The volume is now fixed, the vacuum is supersymmetric and anti-de Sitter, and
$W_0\neq0$ is \emph{necessary}: the $(0,3)$ flux that broke supersymmetry in the no-scale model
is what balances the exponential.

\begin{example}[How small must $W_0$ be?]\label{ex:stabilization:kkltnum}
For $B=1$ and $W_0=-10^{-4}$, solving \eqref{eq:stabilization:kkltW0} numerically gives
$\sigma_{\rm cr}\approx1.81$ and $V_{\min}\approx-5.0\times10^{-10}$. This is not a large
volume: the supergravity approximation needs $\sigma\gg1$. Because $\sigma_{\rm cr}$ grows
only like $\frac1{2\pi}\ln(1/|W_0|)$, a large volume needs an exponentially small $|W_0|$, or
a smaller exponent. KKLT's illustrative figure replaces $2\pi$ by $a=0.1$ (as for gaugino
condensation in a gauge group of large rank); \source{Gra \S7.4, ll.~5027--5029} then
$W_0=-B\ee^{-a\sigma}(1+\tfrac23a\sigma)$, and the same input gives $\sigma_{\rm cr}\approx113.6$
and $V_{\min}=-a^2B^2\ee^{-2a\sigma_{\rm cr}}/(6\sigma_{\rm cr})\approx-2.0\times10^{-15}$
(numerical root finding with mpmath).
\end{example}
Gra\~na stresses that this construction ``should be taken only as a toy model of complete
moduli stabilization'' and reviews the critiques. \source{Gra \S7.3, ll.~4919--4926} Whether
the further step to a de Sitter vacuum is consistent is the subject of the de Sitter
conjecture of \cref{ch:swampland}.

\section{The dual-scale toy potential}\label{sec:stabilization:dualscale}

\subsection{Definition and elementary properties}

None of the potentials above knows about T-duality: the flux potential runs away at large
volume, and at small volume the supergravity approximation fails. This programme studies, as
a model of what a T-duality invariant potential for a radius could look like, the function
\begin{equation}\label{eq:stabilization:dsV}
 V_{\rm ds}(R)=V_0\Bigl(\frac{R}{\sqrt{\ap}}-\frac{\sqrt{\ap}}{R}\Bigr)^{2}
 =\frac{V_0}{\ap}\Bigl[\Bigl(R+\frac{\ap}{R}\Bigr)^2-4\ap\Bigr],\qquad V_0>0,\ R>0 .
\end{equation}
The second form shows that $V_{\rm ds}$ is a function of the dual-scale variable
$R_{\rm eff}=R+\ap/R$ of \cref{ch:dualscale} alone, namely
$V_{\rm ds}=V_0\,(R_{\rm eff}^2-R_{\rm eff,min}^2)/\ap$ with $R_{\rm eff,min}=2\sqrt{\ap}$. Its
properties are calculus:\index{dual-scale potential}\index{self-dual radius}
\begin{enumerate}[leftmargin=*,itemsep=1pt]
\item \emph{Duality invariance.} $V_{\rm ds}(\ap/R)=V_{\rm ds}(R)$, because $R\mapsto\ap/R$
 changes the sign of the bracket in the first form.
\item \emph{Non-negativity and unique minimum.} $V_{\rm ds}\ge0$ with equality if and only if
 $R^2=\ap$, that is, at the self-dual radius $R=\sqrt{\ap}$. Equivalently
 $R+\ap/R\ge2\sqrt{\ap}$, since $R+\ap/R-2\sqrt{\ap}=(R-\sqrt{\ap})^2/R$.
\item \emph{Logarithmic variable.} With $x=\ln(R/\sqrt{\ap})$, T-duality is $x\mapsto-x$ and
 $V_{\rm ds}=4V_0\sinh^2x$, an even function, equal to $4V_0x^2(1+O(x^2))$ near the minimum. The
 variable $x$ has a canonical kinetic term up to a constant, since the kinetic term of
 a radius is proportional to $(\partial R)^2/R^2=(\partial x)^2$ (\cref{ch:backgrounds}), so
 the mass squared of the radion is proportional to $V''(0)=8V_0$ and is positive.
\item \emph{No runaway.} $V_{\rm ds}\to\infty$ both for $R\to\infty$ and for $R\to0$; this is
 what T-duality enforces and what the flux potential of \cref{fig:stabilization:toy} lacks.
\end{enumerate}
\Cref{fig:stabilization:ds} shows the potential. For a $d$-torus with metric $G$ (in units
$\ap=1$) and $B=0$ the natural generalization is
$V_{\rm ds}(G)=V_0\,(\Tr G+\Tr G^{-1}-2d)=V_0\,\Tr\bigl[(G-1)G^{-1}(G-1)\bigr]$, which is
invariant under $G\mapsto G^{-1}$, non-negative for positive definite $G$, and zero at
$G=1$; for $d=1$, $G=R^2$ it reduces to \eqref{eq:stabilization:dsV}. The quantity
$\Tr G+\Tr G^{-1}$ is the trace of the generalized metric $\HH(G,0)$ of \cref{ch:genmetric}.

\begin{figure}[t]
\centering
\begin{tikzpicture}[scale=1.0]
 \draw[-{Stealth}] (-3.4,0) -- (3.6,0) node[below] {$x=\ln(R/\sqrt{\ap})$};
 \draw[-{Stealth}] (0,0) -- (0,3.7) node[right] {$V_{\rm ds}/4V_0$};
 \draw[thick,domain=-1.35:1.35,samples=80] plot ({2*\x},{sinh(\x)*sinh(\x)});
 \draw[dashed,domain=-1.8:1.8,samples=40] plot ({2*\x},{\x*\x});
 \fill (0,0) circle (1.6pt);
 \node[below] at (0,-0.05) {\small $R=\sqrt{\ap}$};
 \draw (2,0.06)--(2,-0.06) node[below] {\small $1$};
 \draw (-2,0.06)--(-2,-0.06) node[below] {\small $-1$};
 \draw[{Stealth}-{Stealth}] (-1.6,2.6) -- (1.6,2.6);
 \node[above] at (0,2.6) {\small $R\leftrightarrow\ap/R$};
 \node at (-2.9,0.6) {\small winding};
 \node at (-2.9,0.25) {\small light};
 \node at (3.0,0.6) {\small momentum};
 \node at (3.0,0.25) {\small light};
\end{tikzpicture}
\caption{The dual-scale toy potential $V_{\rm ds}=4V_0\sinh^2x$ (solid) and its quadratic
approximation $4V_0x^2$ (dashed). T-duality is the reflection $x\mapsto-x$; the unique minimum
is its fixed point, the self-dual radius.}\label{fig:stabilization:ds}
\end{figure}

\subsection{Where could such a potential come from?}

There is a simple observation behind \eqref{eq:stabilization:dsV}. In the conventions of
\cref{ch:circle} a state with momentum number $n$ and winding number $w$ contributes
$n^2/R^2+w^2R^2/\ap^2$ to $M^2$. For $n^2=w^2=1$,
\begin{equation}\label{eq:stabilization:gas}
 \ap\Bigl(\frac{1}{R^2}+\frac{R^2}{\ap^2}\Bigr)-2=\Bigl(\frac{R}{\sqrt{\ap}}-
 \frac{\sqrt{\ap}}{R}\Bigr)^2=\frac{V_{\rm ds}(R)}{V_0} .
\end{equation}
A population with equal numbers of momentum and winding quanta has a mass-squared budget that
is minimized at the self-dual radius: momentum modes are cheap at large $R$, winding modes at
small $R$, and the compromise is $R=\sqrt{\ap}$. This is the mechanism of string gas cosmology\index{string gas cosmology}:
the negative pressure of winding modes drives the radius to smaller values, the positive
pressure of momentum modes to larger values, and ``for an equal number of winding and momentum
modes [\dots] the evolution is driven to the critical radius, the so-called self-dual radius
[\dots] where the total pressure vanishes and T-duality is restored''.
\source{BW \S III.C, ll.~1109--1122} \TL{} The same authors point out what the mechanism does
\emph{not} do: it operates in the string frame and needs a rolling dilaton, which is
``problematic'' at late times, and the dilaton ``still remains a serious challenge''.
\source{BW \S III.D, ll.~1183--1190; \S V, ll.~1957--1958} \Cref{ch:cosmology} returns to this.

We therefore classify as follows. The statements 1--4 above are mathematics, \TA{} where a
Lean declaration is named in \cref{sec:stabilization:lean} and elementary otherwise. That a
gas of momentum and winding modes favours the self-dual radius is \TL. That
\eqref{eq:stabilization:dsV} is the effective potential of a radius in a compactification on
$\KT$, in the Einstein frame, with the dilaton fixed, is \TC: it has not been derived from
\eqref{eq:stabilization:V} with any $K$ and $W$, and three obstacles are known. (i) With
sixteen supercharges no potential is allowed (\cref{sec:stabilization:problem}); the
programme's own Lean documentation records that the scenario is ``in tension with $N=4$
non-renormalization''. (ii) At the self-dual radius additional states become massless
(\cref{ch:tduality}), so an effective description in terms of $R$ alone breaks down exactly
at the minimum; BW stress that these states ``should therefore be considered in the low energy
action''. \source{BW \S III.C, ll.~1135--1138} (iii) The minimum has $V_{\rm ds}=0$, a Minkowski
vacuum: by itself it gives no dark energy (\cref{ch:swampland}).

\section{What the machine has checked}\label{sec:stabilization:lean}

Four Lean modules touch this chapter. They are very different in strength, and the docstrings
of the files themselves say so; we follow them. Proof assistants are introduced in
\cref{ch:leanprimer}.

\subsection{The frontier file on the F-term potential}

The module \lean{StringTheoryFormalization/Frontier/FTermPotential.lean} defines a
structure \lean{DilatonAxion} (a complex $\tau$ with $0<\operatorname{Im}\tau$), a real
function \lean{kahlerPotential} written as $-\log\bigl(-\operatorname{Im}(\tau-\bar\tau)\bigr)$, and
the following toy superpotential. (A sign slip is worth noting: $\operatorname{Im}(\tau-\bar\tau)=
2\operatorname{Im}\tau$, so the argument of the logarithm is negative. Mathlib's \lean{Real.log}
is defined through the absolute value, so the value is nevertheless $-\ln(2\operatorname{Im}\tau)$,
which is the $\tau$ term of \eqref{eq:stabilization:Ktot}; the intended expression is
$-\ln[-\ii(\tau-\bar\tau)]$ with $-\ii(\tau-\bar\tau)=2\operatorname{Im}\tau$ real and positive.)
\begin{leancode}
structure GVWSuperpotential where
  f3 : ℤ
  h3 : ℤ
  period : ℂ
  complexMod : ℂ

noncomputable def GVWSuperpotential.eval (W : GVWSuperpotential) (τ : ℂ) : ℂ :=
  (W.f3 : ℂ) * W.period - τ * (W.h3 : ℂ) * W.period

noncomputable def fTermCondition (W : GVWSuperpotential) (τ : ℂ) : ℂ :=
  -(W.h3 : ℂ) * W.period
\end{leancode}
Thus \lean{eval} is $W(\tau)=(f_3-\tau h_3)\Pi$, the one-period case of
\eqref{eq:stabilization:toyW}, and \lean{fTermCondition} is $\partial_\tau W=-h_3\Pi$: the
plain derivative, \emph{without} the term $(\partial_\tau K)W$ of \eqref{eq:stabilization:V}.
It does not depend on $\tau$, and \lean{kahlerPotential} is not used by any theorem. The three
theorems named in the brief of this chapter are:
\begin{leancode}
theorem fterm_potential_nonneg (W : GVWSuperpotential) (τ : ℂ) (hτ : 0 < τ.im) :
    (0 : ℝ) ≤ ‖fTermCondition W τ‖ ^ 2 := by ...

theorem no_scale_identity :
    (3 : ℚ) = 3 := rfl

theorem fterm_potential_minimum_susy
    (W : GVWSuperpotential) (τ : ℂ)
    (hW : W.eval τ = 0) (hDW : fTermCondition W τ = 0) :
    ‖fTermCondition W τ‖ ^ 2 = 0 := by ...
\end{leancode}
\begin{leanbox}[Plain-language reading]
\lean{fterm\_potential\_nonneg}: the squared norm of the complex number $-h_3\Pi$ is
non-negative. This is true of every complex number; the proof is the tactic
\lean{positivity}, and the hypothesis on $\operatorname{Im}\tau$ is not used. No $\ee^K$, no
inverse K\"ahler metric and no potential $V$ occur in the statement.

\lean{no\_scale\_identity}: the rational number $3$ equals $3$, by \lean{rfl}. The file's own
docstring calls it ``vacuous'': the identity \eqref{eq:stabilization:noscale} appears only in a
comment. Nothing about a K\"ahler metric is proved.

\lean{fterm\_potential\_minimum\_susy}: \emph{if} $-h_3\Pi=0$, then its squared norm is $0$.
The proof rewrites with the hypothesis \lean{hDW} and calls \lean{simp}; the hypothesis
\lean{hW} is not used. It is one direction only and it does not mention a minimum.

A fourth theorem in the file, \lean{flux\_tadpole\_quantization\_exact}, re-exports the
D7/O7 charge cancellation of \cref{ch:tadpole} and says nothing about fluxes.
\end{leanbox}
These theorems are kernel-checked, so they are \TA{} \emph{for what they literally state}; what
they state is far weaker than \cref{prop:stabilization:crit} or \eqref{eq:stabilization:Vns}.
The mathematics of \cref{sec:stabilization:fterm,sec:stabilization:noscale} is \TL{} and
derived on the page; it is not formalized. A faithful formalization would need: the K\"ahler
potential as a real function on the product of upper half planes; Wirtinger derivatives
$\partial_\rho,\partial_{\bar\rho}$ (or, more simply, real partial derivatives in
$\operatorname{Re}\rho,\operatorname{Im}\rho$); the statement
$K^{\rho\bar\rho}K_\rho K_{\bar\rho}=3$ as a theorem about \lean{Real.log}; the definition of
$V$ from $K$ and $W$; and then $V\ge0$ as a corollary. A realistic first step is the algebraic
core, which needs no analysis at all: with $u=\rho-\bar\rho\neq0$,
$\bigl(-u^2/3\bigr)\cdot\bigl(-3/u\bigr)\cdot\bigl(3/u\bigr)=3$
(\cref{ex:stabilization:leannoscale}). \Cref{ch:open} lists this among the open
formalization tasks.

\subsection{The validation use case: integer shadows}

The module \lean{DualScaleValidation/UseCase1\_ModuliStabilization.lean} works with natural
numbers and integers, in units $\ap=1$ (\lean{alpha\_prime : Nat := 1}).
\begin{leancode}
def effective_dual_scale_numerator (R : Nat) : Nat :=
  R * R + alpha_prime

theorem self_dual_is_global_minimum (R : Nat) (h : R ≥ 1) :
    effective_dual_scale_numerator R ≥ 2 := by ...

def moduli_potential (phi phi_0 : Int) : Int :=
  (phi - phi_0) * (phi - phi_0)

theorem moduli_potential_nonneg (phi phi_0 : Int) :
    moduli_potential phi phi_0 ≥ 0 := by ...

theorem moduli_potential_zero_iff (phi phi_0 : Int) :
    moduli_potential phi phi_0 = 0 ↔ phi = phi_0 := by ...
\end{leancode}
\begin{leanbox}[Plain-language reading]
\lean{self\_dual\_is\_global\_minimum}: for every natural number $R\ge1$, $R^2+1\ge2$. The
proof obtains $R\cdot R\ge1$ from \lean{Nat.mul\_pos} and finishes with \lean{omega}. The
quantity $R^2+1$ is $R\cdot R_{\rm eff}(R)$, the \emph{numerator} of $R+1/R$, not
$R_{\rm eff}$ itself. The inequality $R+1/R\ge2$ is equivalent to $R^2+1\ge2R$, which is a
different (and for $R\ge2$ stronger) statement. The theorem is therefore an arithmetic shadow:
it records the value $2$ at $R=1$ and that the numerator does not drop below it on the
positive integers, where no radius $R<1$ exists and T-duality cannot be expressed.

\lean{moduli\_potential\_nonneg} and \lean{moduli\_potential\_zero\_iff}: for integers,
$(\varphi-\varphi_0)^2\ge0$, with equality if and only if $\varphi=\varphi_0$. This is the
integer shadow of a quadratic potential with a unique minimum, such as the leading term
$4V_0x^2$ of $V_{\rm ds}$. The file documents that an earlier version over \lean{Nat}, with
truncated subtraction, had a whole ray of zeros; the present version over \lean{Int} proves
uniqueness. The module also contains \lean{buscher\_log\_involution} ($-(-x)=x$ for integers)
and \lean{moduli\_vacuum\_stability} (the ``if'' direction of the equivalence above).
\end{leanbox}

\subsection{The real statement: the trace bound}

The module \lean{DualScaleStream2/DualScale/TraceBound.lean}, discussed at length in
\cref{ch:dualscale}, proves over $\R$ the faithful version of property 2 of
\cref{sec:stabilization:dualscale}:
\begin{leancode}
theorem circle_effective_scale_ge_two (R : ℝ) (hR : 0 < R) :
    2 ≤ R + R⁻¹ := by ...

theorem dualScale_ge (G : Matrix (Fin d) (Fin d) ℝ) (hG : G.PosDef) :
    (2 * d : ℝ) ≤ dualScale G := by ...

theorem dualScale_inv (G : Matrix (Fin d) (Fin d) ℝ) (hG : IsUnit G.det) :
    dualScale G⁻¹ = dualScale G := by ...
\end{leancode}
Here \lean{dualScale G} is defined as the trace of the generalized metric at $B=0$ and
\lean{dualScale\_eq} identifies it with $\Tr G+\Tr G^{-1}$. In words: for every positive real
$R$, $R+1/R\ge2$; for every positive definite real $d\times d$ matrix,
$\Tr G+\Tr G^{-1}\ge2d$; and the dual scale is invariant under $G\mapsto G^{-1}$. The proof of
the circle case uses the identity $R+R^{-1}-2=(R-1)^2/R$; the matrix case expands
$\Tr[(G-1)G^{-1}(G-1)]\ge0$ and needs no eigenvalues. Together with \lean{dualScale\_one}
($\mathcal D(1)=2d$) these give: $V_{\rm ds}(G)\ge0$, $V_{\rm ds}(G^{-1})=V_{\rm ds}(G)$ and
$V_{\rm ds}(1)=0$. That $G=1$ is the \emph{only} zero is elementary (\cref{ex:stabilization:unique})
but is not among the compiled theorems.

\subsection{The observables module}

\lean{DualScaleValidation/Observables.lean} imports the use case above and records four
numerical relations as exact arithmetic, for instance
\begin{leancode}
theorem dark_energy_cosmological_constant :
    dark_energy_w0 = -1 ∧ dark_energy_wa = 0 := by ...
\end{leancode}
where \lean{dark\_energy\_w0 : Int := -1} and \lean{dark\_energy\_wa : Int := 0} are
definitions; the theorem unfolds them. The module's docstring classifies every one of these
relations as Tier C, states that the module ``is not a derivation of any physical observable
from first principles'', and adds for this one that an exact minimum with $V(\varphi_0)=0$ ``is
a Minkowski vacuum, not de Sitter'', so that it is ``not by itself a prediction of
$w_0=-1$'' with non-zero dark energy density, and that current data are in tension with
$(w_0,w_a)=(-1,0)$. The conjunction \lean{observables\_windows\_master\_contract} states that
four stored numbers lie in four illustrative windows. We mention the module because the brief
of this chapter lists it; it contains no statement about moduli stabilization beyond what
has been said.

\begin{table}[ht]
\centering\small
\begin{tabular}{@{}p{0.47\textwidth}p{0.11\textwidth}p{0.34\textwidth}@{}}
\toprule
Statement & Tier & Evidence\\
\midrule
$W=\int\Omega\wedge G$ reproduces $G^{0,4}=G^{1,3}=0$ & \TL & GVW \S2.4\\
$V=\ee^K(|DW|^2-3|W|^2)$; $DW=0\Rightarrow\partial V=0$ & \TL & Gra (5.21);
 \cref{prop:stabilization:crit}\\
No-scale identity \eqref{eq:stabilization:noscale}, $V\ge0$, minima at $D_iW=0$ & \TL &
 GKP \S2.3; symbolic check\\
Toy model \eqref{eq:stabilization:tau0}--\eqref{eq:stabilization:W0} & derived & symbolic check\\
KKLT formulas \eqref{eq:stabilization:kkltW0}--\eqref{eq:stabilization:kkltV} & \TL & Gra
 (7.8)--(7.9); symbolic check\\
$\|{-h_3\Pi}\|^2\ge0$; $(3:\Q)=3$ & \TA & \lean{fterm\_potential\_nonneg},
 \lean{no\_scale\_identity}\\
$R\in\N$, $R\ge1\Rightarrow R^2+1\ge2$ & \TA & \lean{self\_dual\_is\_global\_minimum}\\
$(\varphi-\varphi_0)^2=0\iff\varphi=\varphi_0$ over $\Z$ & \TA &
 \lean{moduli\_potential\_zero\_iff}\\
$R+1/R\ge2$; $\Tr G+\Tr G^{-1}\ge2d$ over $\R$ & \TA & \lean{circle\_effective\_scale\_ge\_two},
 \lean{dualScale\_ge}\\
String gas drives $R$ to $\sqrt{\ap}$ (string frame) & \TL & BW \S III.C\\
$V_{\rm ds}$ is the radion potential on $\KT$ & \TC & none\\
\bottomrule
\end{tabular}
\caption{Epistemic status of the statements of this chapter. All Lean theorems listed depend
only on the axioms \lean{propext}, \lean{Classical.choice} and \lean{Quot.sound}. The
identification of any Lean object with a physical system is never Tier A.}
\label{tab:stabilization:tiers}
\end{table}

\begin{summarybox}
\begin{itemize}[leftmargin=*,itemsep=1pt]
\item Moduli are massless scalars; they mediate composition-dependent long-range forces and
 must acquire a potential with an isolated minimum.
\item Quantized flux is a fixed integral class whose Hodge type depends on the moduli; the
 supersymmetry conditions on its type are equations for the moduli. They follow from
 $W=\int\Omega\wedge G$ (GVW), in type IIB $W=\int\Omega\wedge(F_3-\tau H_3)$.
\item The $N=1$ potential is $V=\ee^K(K^{I\bar J}D_IW\overline{D_JW}-3|W|^2)$. Points with
 $D_IW=0$ are critical points with $V=-3\ee^K|W|^2\le0$.
\item For $K=-3\ln[-\ii(\rho-\bar\rho)]$ and $W$ independent of $\rho$,
 $K^{\rho\bar\rho}K_\rho K_{\bar\rho}=3$ cancels $-3|W|^2$: $V\ge0$, the minima are at
 $D_iW=0$ ($i\ne\rho$), the volume is flat, and supersymmetry is broken if $W\ne0$. In ten
 dimensions: $G_3$ imaginary self-dual.
\item In the two-period toy model $\bar\tau_0=A/B$ and $\operatorname{Im}\tau_0>0$ requires
 $N_{\rm flux}>0$; one period is not enough.
\item Fixing the volume needs non-perturbative terms (KKLT), giving a supersymmetric AdS
 vacuum and requiring small $|W_0|$.
\item The dual-scale toy potential $V_0(R/\sqrt{\ap}-\sqrt{\ap}/R)^2=4V_0\sinh^2x$ is
 T-duality invariant with its unique minimum at the self-dual radius. The inequality behind
 it is machine-checked over $\R$; its use as a physical radion potential is a proposal.
\item The Lean theorems \lean{fterm\_potential\_nonneg}, \lean{no\_scale\_identity} and
 \lean{fterm\_potential\_minimum\_susy} are trivialities about a toy; the no-scale identity
 itself is not formalized.
\end{itemize}
\end{summarybox}

\section*{Exercises}
\addcontentsline{toc}{section}{Exercises}

\begin{exercise}[$\star$]
Verify \eqref{eq:stabilization:Krho} and \eqref{eq:stabilization:noscale} by writing
$\rho=a+\ii b$ and $K=-3\ln(2b)$. Repeat with $-3$ replaced by $-p$ and show that
$V=\ee^K\bigl(K^{i\bar\jmath}D_iW\overline{D_jW}+(p-3)|W|^2\bigr)$. For which $p$ can $V$ be
negative?
\end{exercise}

\begin{exercise}[$\star$]
Show that $V_{\rm ds}(R)=4V_0\sinh^2x$ with $x=\ln(R/\sqrt{\ap})$, and compute $V''(x)$ at the
minimum. Show that the three expressions in \eqref{eq:stabilization:gas} agree.
\end{exercise}

\begin{exercise}[$\star\star$]\label{ex:stabilization:sl2}
In the toy model \eqref{eq:stabilization:toyW}, the $\SLZ$ duality of type IIB acts by
$\tau\mapsto(a\tau+b)/(c\tau+d)$ and on the flux vectors by
$\binom{F_3}{H_3}\mapsto\bigl(\begin{smallmatrix}a&b\\c&d\end{smallmatrix}\bigr)\binom{F_3}{H_3}$.
(a) Show that $N_{\rm flux}$ is invariant. (b) Show that the vacuum
$\tau_0$ transforms covariantly, i.e.\ the vacuum of the transformed fluxes is the transform of
$\tau_0$. (c) Map the vacuum $\tau_0=\tfrac12(1+\ii)$ of the text to the fundamental domain and
give the transformed fluxes and the string coupling.
\end{exercise}

\begin{exercise}[$\star\star$]
For the toy model with $A=2$, $B=\ii$ expand \eqref{eq:stabilization:toyV} to second order in
$\delta\tau=\tau-2\ii$ at fixed $\rho$. Using the kinetic term
$-K_{\tau\bar\tau}|\partial\tau|^2$ with $K_{\tau\bar\tau}=1/(4(\operatorname{Im}\tau)^2)$, find
the mass of the canonically normalized axio-dilaton as a function of $\operatorname{Im}\rho$.
What happens to this mass along the flat direction $\operatorname{Im}\rho\to\infty$?
\end{exercise}

\begin{exercise}[$\star\star$]
Derive \eqref{eq:stabilization:kkltW0} and \eqref{eq:stabilization:kkltV}. Then show that
for $\sigma_{\rm cr}\gg1$ one has $\sigma_{\rm cr}\approx\frac{1}{2\pi}\ln(|B/W_0|)$ up to
logarithmic corrections, and estimate the $|W_0|$ needed for $\sigma_{\rm cr}=10$ with $B=1$.
\end{exercise}

\begin{exercise}[$\star\star$]\label{ex:stabilization:unique}
Let $G$ be a real symmetric positive definite $d\times d$ matrix. Show that
$\Tr[(G-1)G^{-1}(G-1)]=0$ implies $G=1$. (Write $G^{-1}=S^2$ with $S$ symmetric positive
definite and consider $M=S(G-1)$.) Conclude that $V_{\rm ds}(G)$ has a unique zero. State the
result as a Lean theorem in the style of \lean{dualScale\_ge}; you need not prove it.
\end{exercise}

\begin{exercise}[$\star\star$, Lean]\label{ex:stabilization:leannoscale}
State and prove in Lean 4 with Mathlib the algebraic core of the no-scale identity:
for \lean{u : ℂ} with \lean{u ≠ 0}, \lean{(-(u\^{}2)/3) * (-3/u) * (3/u) = 3}. (The tactic
\lean{field\_simp} followed by \lean{ring} suffices.) Explain in two sentences what is still
missing between your theorem and \eqref{eq:stabilization:noscale}.
\end{exercise}

\begin{exercise}[$\star\star$, Lean]
The theorem \lean{self\_dual\_is\_global\_minimum} proves $R^2+1\ge2$ for natural numbers
$R\ge1$. State and prove the inequality that actually corresponds to $R+1/R\ge2$, namely
\lean{2 * R ≤ R * R + 1} for \lean{R : ℤ}. (Hint: \lean{nlinarith [sq\_nonneg (R - 1)]}.)
Why is the hypothesis $R\ge1$ no longer needed?
\end{exercise}

\begin{exercise}[$\star\star\star$]
Prove the following converse to the ``one period is not enough'' pitfall. For
$W=\sum_{k=1}^m(f_k-\tau h_k)\Pi_k$ with fixed periods, a supersymmetric-in-$\tau$ vacuum
($D_\tau W=0$ with $\operatorname{Im}\tau>0$) exists if and only if
$\operatorname{Im}(\bar AB)>0$, where $A=\sum f_k\Pi_k$, $B=\sum h_k\Pi_k$. Express
$\operatorname{Im}(\bar AB)$ through the antisymmetric matrix
$\operatorname{Im}(\bar\Pi_k\Pi_l)$ and the fluxes, and compare with $\int H_3\wedge F_3$.
\end{exercise}

\section*{Further reading}
\addcontentsline{toc}{section}{Further reading}
\begin{itemize}[leftmargin=*,itemsep=1pt]
\item The superpotential: \source{GVW \S2.4, ll.~876--1008} for the fourfold, the domain wall
 argument and the remark on no-scale models; \source{GVW \S2.5, ll.~1068--1100} for the
 F-theory decomposition of $G$.
\item Type IIB flux compactifications: \source{GKP \S2.3, ll.~740--920} for $K$, $W$, the
 no-scale potential and the match with the imaginary self-dual condition;
 \source{GKP \S3, ll.~1000--1160} for the conifold example and the stabilization of $\tau$.
\item Review: \source{Gra \S5.2 and \S5.4, ll.~2980--3000, 3075--3096, 3490--3552} for the
 $N=1$ potential, the no-scale condition and the potential in terms of $G_3$;
 \source{Gra \S6, ll.~3905--3930} for the moduli problem; \source{Gra \S7.3--7.4,
 ll.~4855--4926, 5022--5029} for KKLT and its critiques.
\item String gases and the self-dual radius: \source{BW \S III.C--D, ll.~1105--1190}.
\item Lean sources, each with a docstring that states what is and is not proved:
 \begin{itemize}[leftmargin=*,itemsep=0pt]
 \item \lean{StringTheoryFormalization/Frontier/FTermPotential.lean}
 \item \lean{DualScaleValidation/UseCase1\_ModuliStabilization.lean}
 \item \lean{DualScaleValidation/Observables.lean}
 \item \lean{DualScaleStream2/DualScale/TraceBound.lean}
 \end{itemize}
\end{itemize}
